\uuid{7kAv}
\exo7id{192}
\titre{exo7 192}
\auteur{bodin}
\organisation{exo7}
\datecreate{1998-09-01}
\isIndication{false}
\isCorrection{false}
\chapitre{Injection, surjection, bijection}
\sousChapitre{Injection, surjection}
\cpge{mpsi-mp2i.raisonnement-vocabulaire-ensembliste.c, pcsi.raisonnement-vocabulaire-ensembliste.d, ptsi.raisonnement-vocabulaire-ensembliste.d, mpsi-mp2i.nombres-complexes.e, pcsi.nombres-complexes.e, ptsi.nombres-complexes.e, mpsi-mp2i.nombres-complexes.c, pcsi.nombres-complexes.c, ptsi.nombres-complexes.c, bcpst1.vocabulaire-applications, bcpst1.nombres-complexes.b}
\module{Algèbre}
\niveau{L1}
\difficulte{}

\contenu{
\question{
L'application $ f : \Cc\setminus{\{0\}} \rightarrow \Cc, \ z\mapsto z+1/z$ est-elle injective ?
surjective ? bijective ?  \\
Donner l'image par $f$ du cercle de centre $0$ et de rayon $1$. \\
Donner l'image r\'eciproque par $f$ de la droite $i\Rr$.
}
}
